\section{Bayesian Receiver Design}
\label{sec:algorithm}

We first estimate the phase increments and noise variance from the
received pilots. With these estimates fixed, the receiver uses message
passing (MP) between a linear estimator and a nonlinear denoiser.
The two modules iteratively exchange extrinsic messages \cite{ma2015}. These messages are obtained
by Gaussian projection and message division \cite{minka2001}. The nonlinear denoiser uses the
learned channel moments and the density ratio for activity detection. The corresponding
learning objectives are presented in Section~\ref{sec:learning}.

\subsection{Parameter Estimation}
\label{sec:frontend}

\subsubsection{Inter-frame phase estimation}
Stack the received samples as
$[\mathbf Y_{\rm s}]_{qN+r,m}=Y_{qmr}$, so that
$\mathbf Y_{\rm s}\in\mathbb C^{QN\times M}$.
To estimate the phase increments, we approximate the contribution of each
candidate terminal by one dominant path with an unconstrained complex
gain at each antenna.
Absorbing the common scalar phase into the gain, we define the path matrix as
\begin{equation}
 \mathbf G_0(d,\nu)=e^{-j2\pi c_1d^2}\mathbf G(d,\nu;0),\qquad
 [\mathbf h_k(d,\nu)]_{qN+r}
 =e^{j2\pi\nu qT_{\rm frame}}
 \sum_{n=0}^{N-1}[\mathbf G_0(d,\nu)]_{rn}P_{qkn}.
 \label{eq:phase-atom}
\end{equation}
The delay search is restricted to $d\in[0,N)$ and the Doppler search to $|\nu|\leq\nu_{\max}$. Under $2Nc_1\in\mathbb N_+$, shifting
the delay by an integer multiple of $N$ changes the path operator by
a scalar phase, which the fitted gains absorb. With
$\mathbf h_k=\mathbf h_k(d_k,\nu_k)$ and gain vector
$\mathbf b_k\in\mathbb C^M$, the fitting objective is
\begin{equation}
 \min_{\{d_k,\nu_k,\mathbf b_k\}}
 \left\|\mathbf Y_{\rm s}-\sum_{k=1}^K
 \mathbf h_k\mathbf b_k^{\mathsf T}\right\|_F^2.
 \label{eq:phase-fit}
\end{equation}
When updating terminal $k$, we fix the other contributions and define
$\mathbf Z_k=\mathbf Y_{\rm s}-\sum_{j\ne k}
\mathbf h_j\mathbf b_j^{\mathsf T}$. For a given $(d,\nu)$, the
least-squares gain and the resulting objective are given by
\begin{equation}
 \begin{aligned}
 \mathbf b_k^{\mathsf T}(d,\nu)
 &=\frac{\mathbf h_k(d,\nu)^{\mathsf H}\mathbf Z_k}
 {\|\mathbf h_k(d,\nu)\|^2},\\
 J_k(d,\nu)&=\|\mathbf Z_k\|_F^2-
 \frac{\|\mathbf h_k(d,\nu)^{\mathsf H}\mathbf Z_k\|^2}
 {\|\mathbf h_k(d,\nu)\|^2}.
 \end{aligned}
 \label{eq:phase-variable-projection}
\end{equation}
The gains are eliminated from the separable nonlinear least-squares
problem by variable projection \cite[Secs.~1 and 2]{oleary2013}. Coarse correlations over
delay, within-frame Doppler, and frame phase provide initial candidates.
A phase candidate $\beta$ corresponds to Doppler aliases
\begin{equation}
 \nu_j=\frac{\beta+2\pi j}{2\pi T_{\rm frame}},\qquad
 j\in\mathbb Z,\quad |\nu_j|\leq\nu_{\max}.
 \label{eq:phase-aliases}
\end{equation}
Candidates are ranked using the pilot response vector in
\eqref{eq:phase-atom}, which enforces the same Doppler within and across
frames. The selected candidates undergo bounded local refinement by
L-BFGS-B \cite{zhu1997}, with delay restricted to one differentiable
branch of the wrapped waveform. Each coordinate update is accepted
when it does not increase the residual energy, up to numerical tolerance.
After each sweep, all antenna gains are refitted jointly by least squares.

The receiver uses a fixed number of complete sweeps, retaining the previous
position as an additional refinement starting point after the first sweep.
It sets $\widehat\beta_k=\operatorname{wrap}
(2\pi\widehat\nu_kT_{\rm frame})$ and constructs
$\widehat{\mathbf A}$ by \eqref{eq:effective-pilot}.
This local fit provides phase estimates for all candidate terminals.
The nonlinear denoiser then estimates the full channel tensor in
\eqref{eq:channel}.

\subsubsection{Noise variance estimation}
Let $L=QN>M\geq2$ and form the spatial sample covariance
\begin{equation}
 \widehat{\mathbf R}_y=\frac1L\sum_{q=0}^{Q-1}\sum_{r=0}^{N-1}
 \mathbf y_{qr}\mathbf y_{qr}^{\mathsf H},\qquad
 \mathbf y_{qr}=(Y_{q0r},\ldots,Y_{q,M-1,r})^{\mathsf T}.
 \label{eq:spatial-noise-covariance}
\end{equation}
Write its eigenvalues in descending order as
$\lambda_1\geq\cdots\geq\lambda_M$; values below
$\epsilon_{\rm mach}\lambda_1$ are floored at that level before
taking logarithms. For a candidate spatial signal rank $j$, define
\begin{equation}
 \bar\lambda_j=\frac1{M-j}\sum_{i=j+1}^M\lambda_i,\qquad
 \log g_j=\frac1{M-j}\sum_{i=j+1}^M\log\lambda_i.
 \label{eq:noise-tail-statistics}
\end{equation}
The Wax--Kailath minimum description length (MDL) criterion, as reviewed
in \cite[Sec.~II-A, eqs.~(5)--(7)]{nadakuditi2008}, gives
\begin{equation}
 \begin{aligned}
 \widehat r_{\rm s}
 &=\arg\min_{0\leq j<M}\left\{
 L(M-j)(\log\bar\lambda_j-\log g_j)
 +\frac{j(2M-j)}2\log L\right\},\\
 \widehat\sigma^2&=\bar\lambda_{\widehat r_{\rm s}}.
 \end{aligned}
 \label{eq:noise-mdl}
\end{equation}
Under a lower-rank spatial signal plus white-noise model, the criterion
selects the spatial signal rank and estimates the noise variance from
the remaining eigenvalues. Both $\widehat\sigma^2$ and
$\widehat{\mathbf A}$ remain fixed throughout the MP iterations.

\subsection{Linear Estimation}

Let the incoming Gaussian message have mean
$\mathbf X^{\rm pri}\in\mathbb C^{KN\times MN}$ and covariance
$\gamma\I_{KN}$ per measurement vector, where $\gamma>0$.
For $\widehat\sigma^2>0$, the posterior mean and covariance under
the estimated operator in \eqref{eq:working-posterior} are given by
\begin{align}
 \mathbf X_{\rm L}
 &=\mathbf X^{\rm pri}
 +\gamma\widehat{\mathbf A}^{\mathsf H}
 (\widehat\sigma^2\I_Q+\gamma\widehat{\mathbf A}\widehat{\mathbf A}^{\mathsf H})^{-1}
 (\mathbf Y-\widehat{\mathbf A}\mathbf X^{\rm pri}),
 \label{eq:linear-mean}\\
 \bm\Sigma_{\rm L}
 &=(\gamma^{-1}\I_{KN}+\widehat\sigma^{-2}\widehat{\mathbf A}^{\mathsf H}\widehat{\mathbf A})^{-1},
 \qquad t_{\rm L}=\frac{\tr\bm\Sigma_{\rm L}}{KN}.
 \label{eq:linear-variance}
\end{align}
We approximate $\bm\Sigma_{\rm L}$ by $t_{\rm L}\I_{KN}$, where
$t_{\rm L}$ is its normalized trace. Dividing the resulting Gaussian
approximation by the incoming Gaussian message gives
\begin{equation}
 v_{\rm L}=(t_{\rm L}^{-1}-\gamma^{-1})^{-1},\qquad
 \mathbf R_{\rm L}=v_{\rm L}
 \left(\frac{\mathbf X_{\rm L}}{t_{\rm L}}
       -\frac{\mathbf X^{\rm pri}}{\gamma}\right).
 \label{eq:linear-extrinsic}
\end{equation}
For $\widehat{\mathbf A}\ne\mathbf0$, the eigenvalues in
\eqref{eq:linear-variance} give $0<t_{\rm L}<\gamma$, so
$v_{\rm L}$ is positive.

A thin singular value decomposition (SVD)
$\widehat{\mathbf A}=\mathbf U\diag(s_i)\mathbf V^{\mathsf H}$ gives
\begin{equation}
 f_{\rm L}=\frac{1}{KN}\sum_i
 \frac{\gamma s_i^2}{\widehat\sigma^2+\gamma s_i^2},\qquad
 t_{\rm L}=\gamma(1-f_{\rm L}),\qquad
 v_{\rm L}=\frac{t_{\rm L}}{f_{\rm L}}.
 \label{eq:linear-fraction}
\end{equation}
Then $\mathbf R_{\rm L}=\mathbf X^{\rm pri}
+(\mathbf X_{\rm L}-\mathbf X^{\rm pri})/f_{\rm L}$.
The SVD is reused throughout the block. After reshaping by
\eqref{eq:mmv-indexing} and applying the damping in
Section~\ref{sec:damping}, the nonlinear estimator receives
$(R_k,v)$ for each terminal.

\subsection{Channel and Activity Denoising}

For channel and activity denoising, we use the following additive white
Gaussian noise (AWGN) observation model
\cite[Secs.~III-B and V-B1]{cai2026}:
\begin{equation}
 R_k=\alpha_kH_k+W_k,\qquad
 H_k\sim q_{\rm tr},\qquad
 \vect(W_k)\sim\CN(\mathbf0,v\I_D).
 \label{eq:pseudo-awgn}
\end{equation}
This model generates the training samples and approximates the
iterative receiver inputs. Before damping, the linear extrinsic mean
satisfies
\begin{equation}
 \begin{aligned}
 \mathbf R_{\rm L}-\mathbf X
 &=(\mathbf L\widehat{\mathbf A}-\I)(\mathbf X-\mathbf X^{\rm pri})\\
 &\quad+\mathbf L(\mathbf A-\widehat{\mathbf A})\mathbf X
  +\mathbf L\mathbf W,\\
 \mathbf L&=\frac{\gamma}{f_{\rm L}}\widehat{\mathbf A}^{\mathsf H}
 (\widehat\sigma^2\I_Q+\gamma\widehat{\mathbf A}\widehat{\mathbf A}^{\mathsf H})^{-1}.
 \end{aligned}
 \label{eq:linear-query-error}
\end{equation}
The three terms represent feedback error, error in the observation
matrix due to phase estimation, and filtered thermal noise. Equation~\eqref{eq:pseudo-awgn} approximates their
sum by isotropic Gaussian noise independent of the channel.

The input densities for active and inactive terminals are given by
\begin{equation}
 \begin{aligned}
 q_{1,v}(R)&=\int\CN(\vect(R);\vect(H),v\I_D)q_{\rm tr}(dH),\\
 q_{0,v}(R)&=\CN(\vect(R);\mathbf0,v\I_D).
 \end{aligned}
 \label{eq:query-densities}
\end{equation}
Bayes' rule gives the local joint posterior
\begin{equation}
 q(dH,a\mid R,v)=
 \frac{\rho^a(1-\rho)^{1-a}
 \CN(\vect(R);a\vect(H),v\I_D)q_{\rm tr}(dH)}
 {(1-\rho)q_{0,v}(R)+\rho q_{1,v}(R)}.
 \label{eq:local-posterior}
\end{equation}
Hence the local activity probability is
\begin{equation}
 p^\star(R,v)=
 \frac{\rho q_{1,v}(R)}
 {(1-\rho)q_{0,v}(R)+\rho q_{1,v}(R)}.
 \label{eq:local-activity}
\end{equation}
Let $m(R,v)=\E[H\mid R,v,\alpha=1]$ and
$\bm\Sigma_H(R,v)=\operatorname{Cov}(\vect(H)\mid R,v,\alpha=1)$.
Using the laws of total expectation and covariance, we have
\begin{equation}
 \begin{aligned}
 \E[X\mid R,v]&=p^\star m,\\
 \operatorname{Cov}(\vect(X)\mid R,v)
 &=p^\star\bm\Sigma_H+
 p^\star(1-p^\star)\vect(m)\vect(m)^{\mathsf H}.
 \end{aligned}
 \label{eq:exact-local-moments}
\end{equation}
This activity-weighted conditional mean is the standard Bayesian
denoising construction \cite[Sec.~III-C]{liu2018}.
Its evaluation requires channel moments and the density ratio
$q_{1,v}/q_{0,v}$, supplied by the channel and activity estimators,
respectively.

We approximate these quantities using a channel mean
$F_k=F_\theta(R_k,v)$, a scalar conditional mean-squared error
$\widehat\tau_k=\widehat\tau_\psi(R_k,v)$, and an activity logit
$z_\phi(R_k,v,F_k)$. We average the conditional mean-squared error
estimates over the $K$ terminals and apply a positive lower bound:
\begin{equation}
 \tau_{\rm raw}=\frac1K\sum_{k=1}^K\widehat\tau_k,\qquad
 \tau=\max\{\tau_{\rm raw},\tau_{\min}\},
 \label{eq:pooled-variance}
\end{equation}
where $\tau_{\min}>0$. The resulting scalar $\tau$ is used as the
variance per coordinate in the Gaussian approximation for each terminal.
The training targets
are specified in Section~\ref{sec:learning}; Appendix~\ref{app:second-order-training}
derives this averaging operation by fitting the Gaussian approximations.

For each terminal, retain a Gaussian approximation to the channel-prior
factor, called a site \cite[Sec.~2.1]{vehtari2020}, of the form
$\CN(\vect(H_k);\vect(E_k),c_k\I_D)$.
For an input variance within the training range and $0<\tau<v$, dividing the Gaussian approximation to the conditional posterior
by the input likelihood gives
\begin{equation}
 c_k=(\tau^{-1}-v^{-1})^{-1},\qquad
 E_k=c_k\left(\frac{F_k}{\tau}-\frac{R_k}{v}\right).
 \label{eq:site-update}
\end{equation}
This condition ensures positive site precision; otherwise, the previous
site is retained as described in Section~\ref{sec:schedule}.
Combining the Gaussian site with the current input, the conditional
variance and mean are respectively given by
\begin{equation}
 t_k=\frac{c_kv}{c_k+v},\qquad
 \mu_k=\frac{c_kR_k+vE_k}{c_k+v}.
 \label{eq:conditional-moments}
\end{equation}
For an accepted update, $\mu_k=F_k$ and $t_k=\tau$.
If the update is rejected, these quantities instead combine the retained
site with the current input.

Within the training variance range, the activity probability is
\begin{equation}
 \ell_k=z_\phi(R_k,v,F_k)-\logit\pi_{\rm tr},\qquad
 p_k=\sigmoid(\logit\rho+\ell_k),
 \label{eq:activity-probability}
\end{equation}
where $\pi_{\rm tr}$ is the training class prior,
$\logit a=\log[a/(1-a)]$, and $\sigmoid b=(1+e^{-b})^{-1}$.
Section~\ref{sec:density-ratio} derives this prior correction.
The resulting local posterior approximation can be expressed as
\begin{equation}
 \begin{aligned}
 \widetilde q_k(dX_k)&=(1-p_k)\delta_0(dX_k)
 +p_k\CN(\vect(X_k);\vect(\mu_k),t_k\I_D)\,dX_k,\\
 \widehat X_k&=p_k\mu_k.
 \end{aligned}
 \label{eq:mixture}
\end{equation}
Let $\mathbf U_k\in\mathbb C^{M\times N^2}$ be obtained by reshaping
$\mu_k$ with one row per antenna. Averaging the covariance across its $N^2$ columns gives
\begin{equation}
 \mathbf C_{X,k}=p_kt_k\I_M+
 \frac{p_k(1-p_k)}{N^2}\mathbf U_k\mathbf U_k^{\mathsf H}.
 \label{eq:joint-spatial-covariance}
\end{equation}
The first term represents conditional channel uncertainty; the second
accounts for uncertainty in terminal activity.

\subsection{Extrinsic Message Update}

Averaging the traces of $\mathbf C_{X,k}$ over terminals and dividing
by the number of antennas gives the scalar variance used for feedback:
\begin{equation}
 \overline t_X=\frac{1}{KM}\sum_k\tr\mathbf C_{X,k}
 =\frac1K\sum_k\left[
 p_kt_k+p_k(1-p_k)\frac{\norm{\mu_k}_F^2}{D}\right].
 \label{eq:joint-average-variance}
\end{equation}
Let $b=\clip(\overline t_X/v,\varepsilon,1-\varepsilon)$,
with $0<\varepsilon<1/2$. Dividing the Gaussian with mean
$\widehat X_k$ and variance $bv$ by the incoming Gaussian message gives
\begin{equation}
 \gamma_{\rm out}=\frac{vb}{1-b},\qquad
 X_{k,\rm out}=\frac{\widehat X_k-bR_k}{1-b}.
 \label{eq:nonlinear-extrinsic}
\end{equation}
Clipping bounds the variance ratio and keeps the outgoing precision
positive. When clipping is inactive, $bv=\overline t_X$.
The damped outgoing parameters become the next linear prior.
